\section{训练与推理的全流程对比}

为了帮助读者建立完整的理解，本节以表格形式对比扩散模型和流模型在训练与推理各环节的差异。

\subsection{训练阶段对比}

\begin{table}[H]
\centering
\caption{扩散模型 vs.\ 流模型：训练阶段对比}
\begin{tabular}{@{} p{3.8cm} p{5cm} p{5cm} @{}}
\toprule
\textbf{环节} & \textbf{扩散模型} & \textbf{流模型（Rectified Flow）} \\
\midrule
时间域 & $t \in \{0, \ldots, T\}$ 或 ${[}0, T{]}$ & $t \in {[}0, 1{]}$ \\
前向过程 & $\mathbf{x}_t = \sqrt{\bar\alpha_t}\,\mathbf{x}_0 + \sqrt{1-\bar\alpha_t}\,\boldsymbol{\epsilon}$ & $\mathbf{x}_t = (1-t)\mathbf{x}_0 + t\mathbf{x}_1$ \\
预测目标 & 噪声 $\boldsymbol{\epsilon}$ 或分数 $\nabla\log p_t$ & 速度 $v = \mathbf{x}_1 - \mathbf{x}_0$ \\
训练损失 & $\|\boldsymbol{\epsilon}_\theta(\mathbf{x}_t,t) - \boldsymbol{\epsilon}\|^2$ & $\|v_\theta(\mathbf{x}_t,t) - (\mathbf{x}_1-\mathbf{x}_0)\|^2$ \\
数据需求 & 仅需 $\mathbf{x}_0 \sim p_{\text{data}}$ & 需要 $\mathbf{x}_0 \sim p_0$ 和 $\mathbf{x}_1 \sim p_1$ \\
\bottomrule
\end{tabular}
\end{table}

\subsection{推理阶段对比}

\begin{table}[H]
\centering
\caption{扩散模型 vs.\ 流模型：推理阶段对比}
\begin{tabular}{@{} p{3.8cm} p{5cm} p{5cm} @{}}
\toprule
\textbf{环节} & \textbf{扩散模型} & \textbf{流模型} \\
\midrule
起点 & $\mathbf{x}_T \sim \mathcal{N}(\mathbf{0},\mathbf{I})$ & $\mathbf{x}_0 \sim \mathcal{N}(\mathbf{0},\mathbf{I})$ \\
求解方式 & SDE 或 PF-ODE & ODE \\
每步更新 & $\mathbf{x}_{t-1} = \mathbf{x}_t + f(\mathbf{x}_t,t)\Delta t$ & $\mathbf{x}_{t+\Delta t} = \mathbf{x}_t + v_\theta(\mathbf{x}_t,t)\Delta t$ \\
典型步数 & 20--100 步 & 1--10 步（Reflow 后） \\
高级求解器 & DPM-Solver 等 & Euler / Midpoint 即可 \\
\bottomrule
\end{tabular}
\end{table}

\subsection{本章小结}

流模型在训练和推理两个阶段都提供了更简洁的流程。特别值得注意的是，Rectified Flow 的推理阶段在 Reflow 后可以用极少步数完成，这是其在实际部署中的最大优势。

